\chapter*{ABSTRACT}
\addcontentsline{toc}{chapter}{\MakeUppercase{ABSTRACT}}

This work addresses the issue of ``social debt'' in Systems Engineering, understood as the accumulation of obligations arising from interpersonal tensions, poor communication, and decisions that prioritize speed over team well-being, generating toxic environments and burnout. Despite its impact on software quality, agile frameworks such as Scrum lack systematic mechanisms to manage these human dynamics.

The objective was to define SocialScrum, an adaptation of Scrum that integrates social debt management into the technical workflow, harmonizing productivity with organizational health. The methodology was based on a Design Science Research approach. A systematic literature review was conducted to characterize ``community smells,'' and the model was designed incorporating the ``Social Guide'' role, metrics such as the ``Health Score,'' the ``Social Product Backlog,'' and ``Social Story Points.'' Conceptual validation was performed through expert judgment with six specialists in agile development, evaluating clarity, completeness, suitability, ease of use, and agility using the AgilityPRO framework.

The results indicated that SocialScrum is a clear proposal consistent with agile principles, with potential to foster collaboration without overloading the process. It is concluded that systematic social debt management is feasible, although empirical validation in real-world scenarios is required. Future work includes pilot implementation through Action Research cycles, development of tools to automate social metrics, and training programs for the proposed role.

\textbf{Keywords:} Social Debt, Agile Software Development, Scrum, SocialScrum, Design Science Research.